% GENERATED FILE. Edit canonical lessons, then run npm run generate:books.
% canonical-source-hash: fnv1a64:b29374061e48b463
% canonical-lessons: MW-C163-plan, MW-C163-sayad, MW-C163-pher, MW-W50-pha, MW-R163-plans

\chapter{Plans and Perhaps}
\label{ch:mw-a2-plans-and-perhaps}
\hlchaptermodality{163}

\begin{chapteropening}
\textbf{By the end of this chapter:} \emph{I can say a plan with its day, say perhaps and again, and write the letter \mw{फ} from \mw{फेर}.}

Complete MW-R163-plans: say a plan with its day, with perhaps and again, and write \mw{फ}.
\end{chapteropening}

\section[\texorpdfstring{\mw{परसूं} \mw{हूं} \mw{बाजार} \mw{जासूं।} (parsū̃ hū̃ bājār jāsū̃.)}{परसूं हूं बाजार जासूं। (parsū̃ hū̃ bājār jāsū̃.)}]{\mw{परसूं} \mw{हूं} \mw{बाजार} \mw{जासूं} — a plan}

\label{lesson:MW-C163-plan}



\begin{quote}
Say \emph{I will go to the market tomorrow}. Then \emph{the day after tomorrow}: \textbf{\mw{परसूं}}.
\end{quote}

\begin{grammarlens}[title={Grammar Lens: the time first}]
A plan opens on its time, then says what you will do:

\begin{quote}
\textbf{\mw{परसूं} \mw{हूं} \mw{बाजार} \mw{जासूं।}}
\end{quote}

\begin{quote}
\emph{parsū̃ hū̃ bājār jāsū̃.} — The day after tomorrow I'll go to the market.
\end{quote}

\begin{quote}
\textbf{\mw{आज} \mw{आथण} \mw{हूं} \mw{घरै} \mw{आसूं।}}
\end{quote}

\begin{quote}
\emph{āj āthaṇ hū̃ gharai āsū̃.} — This evening I'll come home.
\end{quote}
\end{grammarlens}

\subsection*{Your turn}
Two plans join with \textbf{\mw{अर}}:

\begin{quote}
\textbf{\mw{काल} \mw{हूं} \mw{दिनूगै} \mw{उठसूं} \mw{अर} \mw{बाजार} \mw{जासूं।}}
\end{quote}

\begin{quote}
\emph{kāl hū̃ dinūgai uṭhsū̃ ar bājār jāsū̃.} — Tomorrow I'll get up early and go to the market.
\end{quote}

Say these aloud:

\begin{itemize}
\raggedright
  \item \textbf{\mw{परसूं} \mw{हूं} \mw{बाजार} \mw{जासूं।}} — \emph{parsū̃ hū̃ bājār jāsū̃.} — The day after tomorrow I'll go to the market.
  \item \textbf{\mw{आज} \mw{आथण} \mw{हूं} \mw{घरै} \mw{आसूं।}} — \emph{āj āthaṇ hū̃ gharai āsū̃.} — This evening I'll come home.
\end{itemize}

\begin{tcolorbox}[breakable,colback=teal!4,colframe=teal!35!black,title={Before you move on}]
Where does the time go in a plan? (First.) What does the verb take? (\textbf{-\mw{सूं}}.)
\end{tcolorbox}
\section[\texorpdfstring{\mw{सायद} (sāyad)}{सायद (sāyad)}]{\mw{सायद} (sāyad) — perhaps}

\label{lesson:MW-C163-sayad}



\begin{quote}
Before the new one: say \emph{now}, \emph{today} and \emph{every day}. (\textbf{\mw{अबार}}, \textbf{\mw{आज}}, \textbf{\mw{रोज}}.)
\end{quote}

\subsection*{You'll want to know: \mw{सायद}}
\textbf{\mw{सायद}} — \emph{sāyad} — \textquotedblleft{}perhaps\textquotedblright{}.

\begin{quote}
\textbf{\mw{सायद} \mw{हूं} \mw{काल} \mw{आसूं।}}
\end{quote}

\begin{quote}
\emph{sāyad hū̃ kāl āsū̃.} — Perhaps I'll come tomorrow.
\end{quote}

\begin{cousinweb}
\textbf{\mw{सायद}} (\emph{sāyad}) is a loan, usually traced to Persian \emph{shāyad}, \textquotedblleft{}it may be\textquotedblright{}. Marwari says \emph{s} where the Persian has \emph{sh}, as it does in \textbf{\mw{हमेसा}}.
\end{cousinweb}

\begin{grammarlens}[title={What you've built}]
\textbf{\mw{सायद}} opens the sentence and leaves the rest as it was: \textbf{\mw{सायद} \mw{परसूं} \mw{हूं} \mw{बाजार} \mw{जासूं।}} (\emph{sāyad parsū̃ hū̃ bājār jāsū̃.})
\end{grammarlens}

\subsection*{Your turn}
\begin{itemize}
\raggedright
  \item \emph{Say it:} \emph{sāyad}
  \item \emph{Say it:} \textbf{\mw{सायद} \mw{हूं} \mw{काल} \mw{कोनी} \mw{आसूं।}} — \emph{sāyad hū̃ kāl konī āsū̃.} — Perhaps I won't come tomorrow.
  \item \emph{Recall:} say \emph{now}, then \emph{sāyad} again
\end{itemize}

\begin{tcolorbox}[breakable,colback=teal!4,colframe=teal!35!black,title={Before you move on}]
What does \textbf{\mw{सायद}} mean? (\textquotedblleft{}Perhaps\textquotedblright{}.) Say it once more.
\end{tcolorbox}
\section[\texorpdfstring{\mw{फेर} (pher)}{फेर (pher)}]{\mw{फेर} (pher) — again}

\label{lesson:MW-C163-pher}



\begin{quote}
Before the new one: say \emph{perhaps I'll come tomorrow}. Then \emph{later}, \textbf{\mw{पछै}}.
\end{quote}

\subsection*{You'll want to know: \mw{फेर}}
\textbf{\mw{फेर}} — \emph{pher} — \textquotedblleft{}again\textquotedblright{}.

\begin{quote}
\textbf{\mw{हूं} \mw{काल} \mw{फेर} \mw{आसूं।}}
\end{quote}

\begin{quote}
\emph{hū̃ kāl pher āsū̃.} — I'll come again tomorrow.
\end{quote}

It opens on a letter you have not written yet: \emph{ph}, a \textbf{\mw{प}} with a breath after it. The next lesson puts it in your hand.

\begin{cousinweb}
\textbf{\mw{फेर}} (\emph{pher}) is built on the old verb for turning. CDIAL gives Hindi and Gujarati \emph{pher}, \textquotedblleft{}again\textquotedblright{}, under a reconstructed \emph{phirati}, \textquotedblleft{}turns\textquotedblright{}, and lists Old Marwari \emph{phiraï}, \textquotedblleft{}turns\textquotedblright{}, in the same entry (CDIAL 9078).
\end{cousinweb}

\begin{grammarlens}[title={What you've built}]
\textbf{\mw{फेर}} sits right before the verb: \textbf{\mw{फेर} \mw{आसूं}}, \textbf{\mw{फेर} \mw{जासूं}}.
\end{grammarlens}

\subsection*{Your turn}
\begin{itemize}
\raggedright
  \item \emph{Say it:} \emph{pher}
  \item \emph{Say it:} \textbf{\mw{राम} \mw{फेर} \mw{आयो।}} — \emph{rām pher āyo.} — Ram came again.
  \item \emph{Recall:} say \emph{later}, then \emph{pher} again
\end{itemize}

\begin{tcolorbox}[breakable,colback=teal!4,colframe=teal!35!black,title={Before you move on}]
What does \textbf{\mw{फेर}} mean? (\textquotedblleft{}Again\textquotedblright{}.) Say it once more.
\end{tcolorbox}
\section[\texorpdfstring{\mw{फ} (pha)}{फ (pha)}]{\mw{फ} (pha) — the letter that opens \mw{फेर}}

\label{lesson:MW-W50-pha}



\begin{quote}
Before the new one: write \textbf{\mw{प}} once, and name it.

Now say \textbf{\mw{फेर}} (\emph{pher}), \textquotedblleft{}again\textquotedblright{}. It opens on a \textbf{\mw{प}} with a breath after it: \emph{ph}. This lesson puts that letter in your hand.
\end{quote}

\subsection*{Script you'll notice: \mw{फ}}
\textbf{\mw{फ}} — \emph{pha}, the breathed \emph{p}.

It opens \textbf{\mw{फेर}}. It stands beside \textbf{\mw{प}} as \textbf{\mw{ध}} stands beside \textbf{\mw{द}}. Its left part is a stem and a lower bowl joined to a central stem, much like \textbf{\mw{प}}; then comes a separate arch on the right, and the headline.

\subsection*{Writing: \mw{फ}}
\hlblockfigure{\detokenize{figures/MW-W50-pha-filmstrip.pdf}}{How \mw{फ} is written, stroke by stroke}

\begin{itemize}
\raggedright
  \item \textbf{1.} descend the left stem, curve right around the lower bowl, rise along the central side to the headline, then retrace down the central stem without lifting
  \item \textbf{2.} lift at the central junction and sweep clockwise over and down through the open right arch
  \item \textbf{3.} lift and draw the top shirorekhā left-to-right
\end{itemize}

\textbf{Pen lifts: 2.}

\begin{quote}
verified three-stroke teaching form; another learner source stages the components differently
\end{quote}

\begin{quote}
Stroke order is one attested teaching order; hands and teachers vary. Source: JackPotte, ‘Devanagari p\textsuperscript{h} \mw{फ}.gif’, strokes 1–3, Wikimedia Commons, 29 March 2009.
\end{quote}

\subsection*{Your turn}
\begin{itemize}
\raggedright
  \item \emph{Look:} at \textbf{\mw{फेर}} and point at \mw{फ} at its start
  \item \emph{Trace it:} \mw{फ} three times, saying \emph{ph} as you finish each one
  \item \emph{Write it:} \mw{प}, then \mw{फ}, from memory
  \item \emph{Read it:} \textbf{\mw{फेर}} aloud once more
\end{itemize}

\begin{tcolorbox}[breakable,colback=teal!4,colframe=teal!35!black,title={Before you move on}]
Which letter is this — \mw{फ}? (\textbf{pha}, the breathed \emph{p}.) Which word did it come from? (\textbf{\mw{फेर}}, \emph{pher}, \textquotedblleft{}again\textquotedblright{}.)
\end{tcolorbox}
\section[(dialogue)]{Plans and perhaps — a review}

\label{lesson:MW-R163-plans}



\begin{quote}
Say \emph{the day after tomorrow I'll go to the market}.
\end{quote}

\subsection*{Your turn}
\begin{itemize}
\raggedright
  \item \emph{Say it:} \emph{this evening I'll come home}
  \item \emph{Say it:} \emph{tomorrow I'll get up early and go to the market}
  \item \emph{Say it:} \emph{perhaps I'll come tomorrow}; \emph{perhaps I won't come tomorrow}
  \item \emph{Say it:} \emph{I'll come again tomorrow}; \emph{Ram came again}
  \item \emph{Read it:} \textbf{\mw{फेर}}, and name its first letter
\end{itemize}

\begin{tcolorbox}[breakable,colback=teal!4,colframe=teal!35!black,title={Before you move on}]
\emph{Perhaps}? (\textbf{\mw{सायद}}.) \emph{Again}? (\textbf{\mw{फेर}}.) \emph{I'll come}? (\textbf{\mw{आसूं}}.)
\end{tcolorbox}
